\documentclass[12pt,oneside]{book}

%============================================================
% PACKAGES
%============================================================

%--------------------
% Page + font
%--------------------
\usepackage[margin=1in]{geometry}
\usepackage[T1]{fontenc}
\usepackage[utf8]{inputenc}
\usepackage{lmodern}
\usepackage{microtype}
\usepackage{parskip}


%--------------------
% Math
%--------------------
\usepackage{amsmath, amssymb, amsfonts, amsthm, mathtools}
\usepackage{bm}
\usepackage{bbm}
\usepackage{mathrsfs}
\usepackage{dsfont}

%--------------------
% Tables + figures
%--------------------
\usepackage{graphicx}
\usepackage{float}
\usepackage{booktabs}
\usepackage{array}
\usepackage{xltabular}
\usepackage{multirow}
\usepackage{caption}
\usepackage{subcaption}

%--------------------
% Lists
%--------------------
\usepackage{enumitem}
\setlist[itemize]{topsep=2pt,itemsep=2pt}
\setlist[enumerate]{topsep=2pt,itemsep=2pt}

%--------------------
% Colors + hyperlinks
%--------------------
\usepackage[dvipsnames]{xcolor}
\usepackage[
    colorlinks=true,
    linkcolor=MidnightBlue,
    citecolor=MidnightBlue,
    urlcolor=MidnightBlue
]{hyperref}
\usepackage[nameinlink,capitalise]{cleveref}

%--------------------
% Algorithms
%--------------------
\usepackage{algorithm}
\usepackage{algpseudocode}

%--------------------
% Code blocks
%--------------------
\usepackage{listings}

\lstdefinestyle{codestyle}{
    basicstyle=\ttfamily\small,
    breaklines=true,
    frame=single,
    numbers=left,
    numberstyle=\tiny,
    keywordstyle=\color{MidnightBlue},
    commentstyle=\color{ForestGreen},
    stringstyle=\color{BrickRed},
    showstringspaces=false,
    tabsize=4
}

\lstset{style=codestyle}

%--------------------
% Section styling
%--------------------
\usepackage{titlesec}
\usepackage{tocloft}

\setcounter{tocdepth}{2}
\setcounter{secnumdepth}{3}

%============================================================
% THEOREMS
%============================================================

\numberwithin{equation}{section}

\theoremstyle{definition}
\newtheorem{definition}{Definition}[section]
\newtheorem{example}{Example}[section]
\newtheorem{exercise}{Exercise}[section]
\newtheorem{remark}{Remark}[section]

\theoremstyle{plain}
\newtheorem{theorem}{Theorem}[section]
\newtheorem{lemma}{Lemma}[section]
\newtheorem{proposition}{Proposition}[section]
\newtheorem{corollary}{Corollary}[section]

\theoremstyle{remark}
\newtheorem*{note}{Note}

%============================================================
% CUSTOM COMMANDS
%============================================================

% Sets
\newcommand{\R}{\mathbb{R}}
\newcommand{\N}{\mathbb{N}}
\newcommand{\Z}{\mathbb{Z}}
\newcommand{\Q}{\mathbb{Q}}
\newcommand{\C}{\mathbb{C}}

% Probability
\newcommand{\E}{\mathbb{E}}
\newcommand{\Var}{\mathrm{Var}}
\newcommand{\Cov}{\mathrm{Cov}}
\newcommand{\Corr}{\mathrm{Corr}}
\newcommand{\Prob}{\mathbb{P}}

\newcommand{\ind}{\mathbbm{1}}
\newcommand{\iid}{\overset{\text{iid}}{\sim}}
\newcommand{\convd}{\overset{d}{\to}}
\newcommand{\convp}{\overset{p}{\to}}
\newcommand{\as}{\overset{a.s.}{\to}}

% Operators
\newcommand{\given}{\,\middle|\,}
\newcommand{\argmin}{\operatorname*{arg\,min}}
\newcommand{\argmax}{\operatorname*{arg\,max}}

% Linear algebra / analysis
\newcommand{\norm}[1]{\left\lVert #1 \right\rVert}
\newcommand{\abs}[1]{\left\lvert #1 \right\rvert}
\newcommand{\inner}[2]{\left\langle #1, #2 \right\rangle}

\newcommand{\dd}{\,\mathrm{d}}
\newcommand{\grad}{\nabla}
\newcommand{\Hess}{\nabla^2}

\newcommand{\tr}{\operatorname{tr}}
\newcommand{\rank}{\operatorname{rank}}
\newcommand{\diag}{\operatorname{diag}}

% Distributions
\newcommand{\Normal}{\mathcal{N}}
\newcommand{\Bern}{\mathrm{Bernoulli}}
\newcommand{\Bin}{\mathrm{Binomial}}
\newcommand{\Pois}{\mathrm{Poisson}}
\newcommand{\Gam}{\mathrm{Gamma}}
\newcommand{\BetaDist}{\mathrm{Beta}}
\newcommand{\Exp}{\mathrm{Exponential}}
\newcommand{\Dir}{\mathrm{Dirichlet}}
\newcommand{\InvGam}{\mathrm{Inv\text{-}Gamma}}

% Title block
\title{EN.553.634 Elements of Statistical Learning}
\author{Jonathan Y. Ma}
\date{June 2026}

\begin{document}

\maketitle

\noindent \emph{This course was taught by Dr. Sergey Kushnarev in the Fall 2026 Semester. The notes are transcribed by Jonathan Ma and may contain errors and omitted content, so any issues are to be directed towards me.}

\tableofcontents

\section{Introduction to Statistical Learning}

Let $(X,Y)\sim P$, where $X\in\mathcal X\subseteq\R^p$ is a feature vector
and $Y$ is a response. Statistical learning uses a sample
$\mathcal D_n=\{(X_i,Y_i)\}_{i=1}^n$ to construct a prediction rule
$f:\mathcal X\to\mathcal A$. In regression, $\mathcal A=\R$; in
classification, $\mathcal A$ is a finite set or a vector of class
probabilities.

Prediction and interpretation are different goals. A predictive model is
judged on new observations from the deployment distribution. A coefficient
can be interpreted causally only with additional identification assumptions;
predictive accuracy alone does not provide them.

\begin{example}[From coin flips to supervised learning]
Suppose a coin's probability of heads depends on measurable conditions $X$,
and $Y\in\{0,1\}$ records the outcome. Estimating one bias $p$ becomes
estimating the conditional probability
\[
    \eta(x)=P(Y=1\mid X=x).
\]
A classifier converts $\eta(x)$ into an action, while a calibrated probability
model attempts to estimate $\eta(x)$ itself.
\end{example}

\subsection{Loss Functions}

A loss $L(y,a)$ quantifies the cost of predicting $a$ when the outcome is
$y$. Common choices are
\begin{align*}
 L_{\mathrm{sq}}(y,a)&=(y-a)^2,&
 L_{\mathrm{abs}}(y,a)&=\abs{y-a},\\
 L_{01}(y,a)&=\ind\{y\neq a\},&
 L_{\mathrm{log}}(y,q)&=-y\log q-(1-y)\log(1-q).
\end{align*}
Squared loss strongly penalizes large residuals; absolute loss is more robust.
Zero-one loss targets classification accuracy but is discontinuous and hard
to optimize. Logistic loss is a smooth convex surrogate and estimates class
probabilities when the model is well specified.

A loss should reflect the decision. For asymmetric binary costs $c_{10}$ for
a false positive and $c_{01}$ for a false negative,
\[
 L(y,a)=c_{10}\ind\{y=0,a=1\}+c_{01}\ind\{y=1,a=0\}.
\]

\subsection{Risk and Empirical Risk Minimization}

The population and empirical risks are
\begin{equation}
 R(f)=\E[L(Y,f(X))],
 \qquad
 \widehat R_n(f)=\frac1n\sum_{i=1}^nL(Y_i,f(X_i)).
 \label{eq:learning-risk}
\end{equation}
Given a function class $\mathcal F$, empirical risk minimization (ERM) chooses
$\widehat f\in\argmin_{f\in\mathcal F}\widehat R_n(f)$. If $f^*$ minimizes
risk over all measurable rules, then
\[
 R(\widehat f)-R(f^*)=
 \underbrace{R(\widehat f)-\inf_{f\in\mathcal F}R(f)}_{\text{estimation error}}
 +\underbrace{\inf_{f\in\mathcal F}R(f)-R(f^*)}_{\text{approximation error}}.
\]
The first term reflects finite data; the second reflects the restricted model
class. Interpolation of training data alone gives no guarantee of
generalization.

\subsection{Optimal Prediction}

Conditioning on $X=x$ reduces prediction to a pointwise decision:
\begin{equation}
    f^*(x)\in\argmin_a\E[L(Y,a)\mid X=x].
    \label{eq:bayes-predictor}
\end{equation}
Under squared loss, $f^*(x)=\E[Y\mid X=x]$; under absolute loss, any
conditional median is optimal. Under zero-one loss, the optimal classifier
chooses the most probable class.

\begin{proof}[Squared-loss calculation]
Writing $m(x)=\E[Y\mid X=x]$,
\[
 \E[(Y-a)^2\mid X=x]
 =\Var(Y\mid X=x)+\{m(x)-a\}^2,
\]
which is minimized by $a=m(x)$.
\end{proof}

\subsection{The Regression Function}

The regression function $m(x)=\E[Y\mid X=x]$ describes the systematic part
of a response. Writing
\[
    Y=m(X)+\varepsilon,\qquad \E[\varepsilon\mid X]=0,
\]
does not require normality or constant variance. The conditional variance
$\sigma^2(x)=\Var(Y\mid X=x)$ may depend on $x$. For Bernoulli $Y$,
$m(x)=P(Y=1\mid X=x)$ is both the regression function and the class
probability.

\subsection{Generalization Error}

Conditional on the training data, test error is
\[
 \operatorname{Err}(\mathcal D_n)
 =\E[L(Y^{\mathrm{new}},\widehat f(X^{\mathrm{new}}))\mid\mathcal D_n].
\]
Training error reuses observations that selected $\widehat f$ and is
optimistically biased. Test observations must be excluded from every modeling
choice, including imputation, scaling, feature selection, and tuning.

Ordinary validation assumes that training and deployment observations have
the same distribution. Covariate shift changes $P_X$ but not $P_{Y\mid X}$;
concept drift changes the conditional law. Random cross-validation on
historical data does not diagnose either automatically.

\subsection{The Bias-Variance Tradeoff}

For squared-error prediction at fixed $x$, averaging over training sets gives
\begin{align}
 \E[(Y_0-\widehat f(x))^2\mid X_0=x]
 &=\sigma^2(x)
 +\{\E[\widehat f(x)]-m(x)\}^2
 +\Var(\widehat f(x)).
 \label{eq:prediction-bias-variance}
\end{align}
These are irreducible noise, squared bias, and variance. Restrictive models
tend toward higher bias and lower variance; flexible models tend toward lower
approximation bias and higher sampling variance. This exact additive
decomposition is specific to squared loss.

\section{Least Squares}

Given $X\in\R^{n\times p}$ and $y\in\R^n$, ordinary least squares solves
\begin{equation}
    \widehat\beta\in\argmin_{\beta\in\R^p}\norm{y-X\beta}^2.
    \label{eq:ols-objective}
\end{equation}
The normal equations are
$X^\mathsf TX\widehat\beta=X^\mathsf Ty$. If $X$ has full column rank,
\[
    \widehat\beta=(X^\mathsf TX)^{-1}X^\mathsf Ty.
\]
Computationally, solve the system using QR or SVD rather than explicitly
forming the inverse; forming the Gram matrix squares the condition number.

\subsection{Estimation and Geometry}

The fitted vector is the orthogonal projection of $y$ onto the column space:
\[
 \widehat y=Hy,\qquad
 H=X(X^\mathsf TX)^{-1}X^\mathsf T.
\]
The hat matrix is symmetric and idempotent. The residual $e=y-\widehat y$ is
orthogonal to every column of $X$. If $X$ is rank deficient, fitted values
remain unique but coefficients do not; $X^+y$ is the minimum-norm solution.

Leverage $h_{ii}=H_{ii}$ measures how unusual row $i$ is in feature space.
Since $\sum_i h_{ii}=\tr(H)=p$, average leverage is $p/n$. The leave-one-out
residual satisfies
\begin{equation}
    y_i-\widehat y_i^{(-i)}=\frac{e_i}{1-h_{ii}},
    \label{eq:loo-ols}
\end{equation}
so OLS leave-one-out error can be computed without $n$ refits.

\subsection{Sampling Properties of Least Squares}

Under the fixed-design model
\[
 y=X\beta+\varepsilon,\qquad
 \E[\varepsilon\mid X]=0,\qquad
 \Cov(\varepsilon\mid X)=\sigma^2I,
\]
OLS is conditionally unbiased and
\[
    \Cov(\widehat\beta\mid X)=\sigma^2(X^\mathsf TX)^{-1}.
\]
Gauss--Markov says OLS has smallest covariance among linear unbiased
estimators; it does not require Gaussian errors. Normality gives exact $t$ and
$F$ inference. An unbiased error-variance estimator is
\[
    \widehat\sigma^2=\frac{\norm{y-X\widehat\beta}^2}{n-p}.
\]
With heteroskedastic errors, use sandwich standard errors or weighted least
squares. If $\E[\varepsilon\mid X]\neq0$, robust standard errors do not remove
coefficient bias.

\begin{example}[Omitted-variable bias]
If $Y=\beta_0+\beta_1X+\beta_2Z+\varepsilon$ but $Z$ is omitted, the population
slope from regressing $Y$ on $X$ is
\[
    \beta_1+\beta_2\frac{\Cov(X,Z)}{\Var(X)}.
\]
Bias requires both a nonzero effect of $Z$ and association between $X$ and
$Z$.
\end{example}

\section{Multiple Regression}

An intercept, numeric predictors, and encoded categorical variables are
columns of a design matrix. With an intercept, a $K$-level categorical
predictor uses $K-1$ reference indicators. The coefficient of $X_j$ is the
change in fitted response per unit change in $X_j$, holding the other
represented features fixed.

The Frisch--Waugh--Lovell theorem says the coefficient of $X_j$ equals the
slope obtained by residualizing both $Y$ and $X_j$ against the remaining
predictors. Multiple regression therefore compares variation in $X_j$ not
linearly explained by the other columns.

\subsection{Interactions and Basis Expansions}

An interaction permits one effect to depend on another:
\[
    m(x_1,x_2)=\beta_0+\beta_1x_1+\beta_2x_2+\beta_{12}x_1x_2.
\]
Then $\partial m/\partial x_1=\beta_1+\beta_{12}x_2$, so $\beta_1$ is the
effect of $x_1$ only at $x_2=0$. Centering can improve interpretation.

A basis expansion uses fixed transformations $h_m$:
\[
    f(x)=\sum_{m=1}^M\beta_mh_m(x).
\]
The model is nonlinear in $x$ but linear in $\beta$. Raw high-degree
polynomials are often unstable; orthogonal polynomials or splines are usually
better conditioned.

\section{Logistic Regression}

For $Y\in\{0,1\}$, logistic regression models
\begin{equation}
 p(x)=P(Y=1\mid X=x),\qquad
 \log\frac{p(x)}{1-p(x)}=\beta_0+x^\mathsf T\beta.
 \label{eq:logistic-model}
\end{equation}
Thus $p(x)=\{1+\exp[-(\beta_0+x^\mathsf T\beta)]\}^{-1}$. A unit increase in
$x_j$ multiplies conditional odds by $e^{\beta_j}$ holding other modeled
features fixed; it does not add a constant amount to probability.

\subsection{Likelihood}

Including the intercept in $x_i$, the Bernoulli log-likelihood is
\begin{equation}
 \ell(\beta)=\sum_{i=1}^n
 [y_ix_i^\mathsf T\beta-\log\{1+e^{x_i^\mathsf T\beta}\}].
 \label{eq:logistic-likelihood}
\end{equation}
If $p_i=p(x_i)$, then
\[
 \grad\ell(\beta)=X^\mathsf T(y-p),\qquad
 \Hess\ell(\beta)=-X^\mathsf TWX,
\]
where $W=\diag\{p_i(1-p_i)\}$. The likelihood is concave.

Under perfect separation, a hyperplane correctly separates all observations
and the likelihood supremum occurs as $\norm\beta\to\infty$; no finite
unpenalized MLE exists. Ridge regularization gives a finite solution.
Numerically, evaluate $\log(1+e^z)$ with a stable softplus routine.

\subsection{Fitting Logistic Regression: Newton and IRLS}

Newton's update is
\begin{equation}
 \beta^{\mathrm{new}}
 =\beta^{\mathrm{old}}+(X^\mathsf TWX)^{-1}X^\mathsf T(y-p).
 \label{eq:logistic-newton}
\end{equation}
With working response $z=X\beta+W^{-1}(y-p)$, this is weighted least squares:
\[
    \beta^{\mathrm{new}}=(X^\mathsf TWX)^{-1}X^\mathsf TWz.
\]
This is iteratively reweighted least squares (IRLS). Implementations solve the
weighted system rather than invert it, and use step halving or line search if
a full Newton step decreases the likelihood. Growing coefficients and fitted
probabilities near zero or one suggest separation.

\section{Bayes Classification}

Let $Y\in\{1,\ldots,K\}$ and $\eta_k(x)=P(Y=k\mid X=x)$. Under zero-one loss,
the Bayes classifier is
\[
    g^*(x)=\argmax_k\eta_k(x),
\]
with Bayes error $R^*=\E[1-\max_k\eta_k(X)]$.

Generative classifiers model priors $\pi_k$ and class densities $f_k$, then
apply
\[
 P(Y=k\mid X=x)=\frac{\pi_kf_k(x)}{\sum_j\pi_jf_j(x)}.
\]
LDA assumes normal class distributions with a shared covariance, producing
linear boundaries. QDA allows separate covariance matrices and quadratic
boundaries, trading lower bias for substantially more parameters.

\subsection{Thresholds}

With false-positive cost $c_{10}$ and false-negative cost $c_{01}$, predict
class $1$ when
\begin{equation}
    \eta(x)>\frac{c_{10}}{c_{10}+c_{01}}.
    \label{eq:cost-threshold}
\end{equation}
The threshold $1/2$ is optimal only for equal costs. Changing a threshold
changes decisions, not ranking or calibration. Select it on validation data
using deployment costs or capacity constraints.

\subsection{ROC and Calibration}

For threshold $t$,
\[
 \operatorname{TPR}(t)=P(\widehat\eta(X)>t\mid Y=1),\qquad
 \operatorname{FPR}(t)=P(\widehat\eta(X)>t\mid Y=0).
\]
The ROC curve plots TPR against FPR. AUC is the probability that a random
positive receives a higher score than a random negative, with half credit for
ties. It assesses ranking, not probability accuracy, and may hide behavior in
the operational low-FPR region.

A score is calibrated when
$P(Y=1\mid\widehat\eta(X)=q)=q$. Reliability plots, log loss, and Brier score
assess calibration. Platt scaling fits a logistic recalibration; isotonic
regression fits a monotone map. Under class imbalance, precision--recall
curves are often more informative because precision depends on prevalence.

\section{Estimating Generalization Error}

The quantity of interest is performance of the entire training procedure,
including tuning. A test set estimates this only if it is untouched until the
procedure is finalized. Repeatedly checking test performance turns the test
set into training information.

\subsection{Validation and CV}

In a validation split, train on one subset, choose hyperparameters on another,
and evaluate once on a test subset. In $K$-fold cross-validation (CV), partition
indices into folds $I_1,\ldots,I_K$ and compute
\[
 \widehat{\operatorname{Err}}_{\mathrm{CV}}
 =\frac1n\sum_{k=1}^K\sum_{i\in I_k}
 L\{Y_i,\widehat f^{(-k)}(X_i)\}.
\]
Five- or ten-fold CV is a common variance--computation compromise. Leave-one-
out CV has nearly maximal training sets but can have high variance and is not
always preferable.

Preprocessing must be refit within each training fold. For grouped,
longitudinal, spatial, or time-ordered data, folds must respect dependence and
deployment direction. Random row splits can leak subjects or future
information.

Nested CV uses inner folds for tuning and outer folds for evaluation. The
one-standard-error rule chooses the simplest model whose estimated error is
within one standard error of the minimum, favoring stability when differences
are mostly noise.

\subsection{Bootstrap and Statistical Uncertainty}

The nonparametric bootstrap samples $n$ training observations with replacement,
refits the complete procedure, and approximates sampling variability from the
replicates. About $1-e^{-1}\approx0.632$ of distinct observations appear in a
large bootstrap sample; the rest are out-of-bag observations.

Bootstrap standard errors estimate variability of parameters or predictions:
\[
 \widehat{\operatorname{se}}_{\mathrm{boot}}(\widehat\theta)
 =\left\{\frac1{B-1}\sum_{b=1}^B
 (\widehat\theta^{*(b)}-\bar\theta^*)^2\right\}^{1/2}.
\]
Percentile, basic, and BCa intervals make different corrections. The ordinary
bootstrap can fail for nonsmooth statistics, parameters on boundaries,
extreme tails, or dependent observations. Use block bootstrap for suitable
dependent sequences.

CV fold errors are correlated because fitted training sets overlap; treating
them as independent observations understates uncertainty. Comparing two
models by paired out-of-fold losses is better than comparing separate averages,
but repeated or nested resampling is needed for honest selection uncertainty.

\section{Model Selection}

Model selection chooses features, transformations, hyperparameters, or an
algorithm. Training fit always improves as a nested model grows, so selection
must penalize complexity or estimate out-of-sample error.

For Gaussian linear models, common criteria are
\[
 \mathrm{AIC}=-2\ell(\widehat\theta)+2d,\qquad
 \mathrm{BIC}=-2\ell(\widehat\theta)+d\log n,
\]
where $d$ is the number of fitted parameters. AIC targets predictive
Kullback--Leibler risk; BIC more strongly favors small models and, under
restrictive conditions, consistently selects a true finite-dimensional model.
They answer different questions.

Searching many models creates selection bias. Standard errors and $p$-values
computed as if the selected model had been fixed in advance are generally too
optimistic.

\subsection{Complexity and Degree of Freedom}

For a linear smoother $\widehat y=Sy$, effective degrees of freedom are
\[
    \operatorname{df}=\tr(S).
\]
OLS has $S=H$ and $\operatorname{df}=p$. For nonlinear procedures, a covariance
definition is
\[
 \operatorname{df}
 =\frac1{\sigma^2}\sum_{i=1}^n\Cov(\widehat y_i,y_i).
\]
It measures sensitivity of fitted values to observed responses, not simply
the number of nonzero parameters.

Mallows' $C_p$ and Stein's unbiased risk estimate correct training error by an
optimism term proportional to degrees of freedom. Complexity can also be
controlled through norm constraints, tree depth, early stopping, margin, or
algorithmic stability.

\section{Ridge Regression}

Ridge regression solves
\begin{equation}
 \widehat\beta_\lambda
 =\argmin_\beta\{\norm{y-X\beta}^2+\lambda\norm\beta^2\},
 \qquad\lambda\geq0,
 \label{eq:ridge-objective}
\end{equation}
giving
\[
    \widehat\beta_\lambda=(X^\mathsf TX+\lambda I)^{-1}X^\mathsf Ty.
\]
The intercept is normally left unpenalized. Predictors must be standardized
because the penalty depends on units. Ridge is also the posterior mode and
mean under a Gaussian prior in the Gaussian linear model.

\subsection{Shrinkage}

Conditional on $X$ under $y=X\beta+\varepsilon$,
\[
 \E[\widehat\beta_\lambda]
 =(X^\mathsf TX+\lambda I)^{-1}X^\mathsf TX\beta,
\]
so ridge is biased. It can nevertheless lower prediction MSE by greatly
reducing variance, especially under collinearity or $p\geq n$. Unlike lasso,
ridge usually retains every feature.

As $\lambda$ increases, coefficient norm and effective degrees of freedom
decrease. Tune $\lambda$ by validation, fitting standardization within each
fold. The constrained form $\norm\beta^2\leq t$ is equivalent for a matching
Lagrange multiplier.

\subsection{SVD}

Let the centered design have thin SVD $X=UDV^\mathsf T$, with singular values
$d_1\geq\cdots\geq d_r>0$. Then
\[
 \widehat\beta_\lambda
 =V\diag\left\{\frac{d_j}{d_j^2+\lambda}\right\}U^\mathsf Ty,
\quad
 \widehat y_\lambda
 =U\diag\left\{\frac{d_j^2}{d_j^2+\lambda}\right\}U^\mathsf Ty.
\]
Ridge shrinks low-variance directions most strongly. Its effective degrees of
freedom are
\[
    \operatorname{df}(\lambda)=\sum_{j=1}^r\frac{d_j^2}{d_j^2+\lambda}.
\]
The SVD expression is central in interviews: it explains stability, shrinkage,
and why ridge remains well-defined when $X^\mathsf TX$ is singular.

\section{Lasso Regression}

The lasso solves
\begin{equation}
 \widehat\beta_\lambda
 =\argmin_\beta\left\{\frac12\norm{y-X\beta}^2
 +\lambda\norm\beta_1\right\}.
 \label{eq:lasso-objective}
\end{equation}
It is convex but nonsmooth at zero. The $\ell_1$ constraint has corners, so
solutions often lie on coordinate axes and are sparse.

\subsection{Sparsity}

Lasso performs shrinkage and variable selection together. With correlated
predictors it may select one arbitrarily, making the selected set unstable.
Ridge tends to distribute weight across the group. Elastic net combines
$\ell_1$ and squared $\ell_2$ penalties and often stabilizes grouped features.

Exact support recovery needs conditions stronger than those needed for good
prediction, including sufficient signal and restrictions on correlations.
A sparse fitted model should not automatically be interpreted as identifying
the true causal variables.

\subsection{KKT}

Let $r=y-X\widehat\beta$. The Karush--Kuhn--Tucker conditions are
\begin{equation}
 X_j^\mathsf Tr=
 \begin{cases}
   \lambda\operatorname{sign}(\widehat\beta_j),
      &\widehat\beta_j\neq0,\\
   z_j,\quad\abs{z_j}\leq\lambda,
      &\widehat\beta_j=0.
 \end{cases}
 \label{eq:lasso-kkt}
\end{equation}
Thus a variable stays at zero when its correlation with the current residual
does not exceed $\lambda$. These conditions certify optimality and support
screening rules.

If $\lambda\geq\norm{X^\mathsf Ty}_\infty$, the all-zero vector is optimal
when the response and columns have been centered.

\subsection{Coordinate Descent}

Holding all other coefficients fixed, define the partial residual
$r_j=y-\sum_{k\neq j}X_k\beta_k$. The coordinate update is
\begin{equation}
 \beta_j\leftarrow
 \frac{S(X_j^\mathsf Tr_j,\lambda)}{\norm{X_j}^2},
 \qquad
 S(z,\lambda)=\operatorname{sign}(z)(\abs z-\lambda)_+.
 \label{eq:lasso-coordinate}
\end{equation}
The soft-thresholding operator sets small correlations exactly to zero.
Cyclic coordinate descent is efficient for sparse matrices. Warm starts along
a decreasing grid of $\lambda$ values greatly reduce computation; convergence
should be checked through KKT residuals or the primal--dual gap.

\section{Regression Splines and Basis Expansions}

A degree-$d$ regression spline is piecewise polynomial with continuity
constraints at knots $\xi_1,\ldots,\xi_K$. A truncated-power basis is
\[
 1,x,\ldots,x^d,(x-\xi_1)_+^d,\ldots,(x-\xi_K)_+^d.
\]
Cubic splines are twice continuously differentiable. B-spline bases span the
same space but have local support and much better numerical conditioning.

With a fixed basis matrix $B$, fitting is linear regression. Flexibility is
controlled by degree, knot locations, number of knots, or a roughness penalty.
Boundary behavior matters: natural cubic splines constrain the function to be
linear beyond the boundary knots and often extrapolate more sensibly.

\subsection{Smoothing Splines}

The cubic smoothing spline minimizes
\begin{equation}
 \sum_{i=1}^n\{y_i-f(x_i)\}^2
 +\lambda\int\{f''(t)\}^2\dd t
 \label{eq:smoothing-spline}
\end{equation}
over functions with square-integrable second derivative. The minimizer is a
natural cubic spline with knots at the distinct observed $x_i$ values. Thus an
apparently infinite-dimensional optimization has a finite-dimensional
solution.

For a suitable basis, the coefficient problem is
\[
 \widehat\beta=(B^\mathsf TB+\lambda\Omega)^{-1}B^\mathsf Ty,
 \qquad
 \Omega_{jk}=\int B_j''(t)B_k''(t)\dd t.
\]
The fitted values are $S_\lambda y$ and
$\operatorname{df}(\lambda)=\tr(S_\lambda)$. As $\lambda\to0$ the fit
interpolates; as $\lambda\to\infty$ it approaches a straight line.
Generalized cross-validation approximates leave-one-out error by
\[
 \operatorname{GCV}(\lambda)
 =\frac{n^{-1}\norm{(I-S_\lambda)y}^2}
        {\{1-\tr(S_\lambda)/n\}^2}.
\]

\subsection{GAMs}

A generalized additive model (GAM) assumes
\[
    g\{\E(Y\mid X=x)\}
    =\beta_0+\sum_{j=1}^pf_j(x_j),
\]
where $g$ is a link function. Each component is represented by a basis and
regularized with a smoothness penalty. Identifiability commonly uses
$\sum_i f_j(x_{ij})=0$.

GAMs capture nonlinear main effects while retaining component plots and
efficient fitting. Backfitting alternately updates each $f_j$ using partial
residuals. Penalized IRLS handles non-Gaussian responses. The additive
assumption excludes interactions unless terms such as tensor-product smooths
are explicitly included.

\section{Decision Trees}

A regression tree partitions feature space into terminal regions
$R_1,\ldots,R_M$ and predicts
\[
    \widehat f(x)=\sum_{m=1}^M\widehat c_m\ind\{x\in R_m\},
    \qquad
    \widehat c_m=\frac1{\abs{R_m}}\sum_{i:x_i\in R_m}y_i.
\]
At each node, CART greedily chooses a feature $j$ and threshold $s$ minimizing
the sum of squared errors in the two child nodes. For classification, common
node impurities are
\[
 Q_{\mathrm{Gini}}=1-\sum_k\widehat p_k^2,\qquad
 Q_{\mathrm{entropy}}=-\sum_k\widehat p_k\log\widehat p_k.
\]

Trees model interactions automatically, need little preprocessing, and are
invariant to monotone transformations of individual features. A single tree
is unstable: small data changes can alter early splits and the whole
partition. Greedy training is computationally feasible but does not find the
globally optimal tree.

Continuous split search can be done after sorting each feature. Missing
values require an explicit strategy, such as imputation, learned default
directions, or surrogate splits. Unrestricted high-cardinality categorical
features can receive an unfairly large chance to reduce training impurity.

\subsection{Pruning}

Grow a sufficiently large tree $T_0$, then choose a subtree by cost-complexity
pruning:
\begin{equation}
    C_\alpha(T)=\sum_{m=1}^{\abs T}
       \sum_{i:x_i\in R_m}(y_i-\widehat c_m)^2
       +\alpha\abs T.
    \label{eq:tree-pruning}
\end{equation}
Weakest-link pruning produces a nested sequence of optimal subtrees as
$\alpha$ increases. Cross-validation chooses $\alpha$; the one-standard-error
rule often yields a smaller, more stable tree.

Stopping growth early can miss a useful split whose benefit appears only
after another split. Growing then pruning separates search from complexity
selection.

\section{Random Forests}

Random forests average many deep randomized trees:
\[
    \widehat f_{\mathrm{RF}}(x)=\frac1B\sum_{b=1}^B\widehat f_b(x).
\]
Each tree uses a bootstrap sample, and each split considers a random subset of
features. Feature subsampling decorrelates trees; averaging then reduces
variance more effectively.

If identically distributed tree predictions have variance $\sigma^2$ and
pairwise correlation $\rho$, the average of $B$ trees has variance
\[
 \rho\sigma^2+\frac{1-\rho}{B}\sigma^2.
\]
Increasing $B$ removes only the second term. Forest performance therefore
depends on both individual tree strength and inter-tree correlation.

Out-of-bag predictions evaluate observation $i$ using only trees whose
bootstrap samples excluded it. They provide a convenient internal error
estimate, though an external test set is still needed after extensive tuning.

\subsection{Bagging}

Bagging fits the same unstable learner to bootstrap samples and averages:
\[
    \widehat f_{\mathrm{bag}}(x)
    =\frac1B\sum_{b=1}^B\widehat f^{*(b)}(x).
\]
For squared loss, averaging reduces variance. It helps trees because they are
unstable, but changes little for an already stable linear estimator.

Permutation importance measures the loss increase after shuffling a feature
in out-of-bag data. It measures predictive reliance, not causal effect, and
can split importance among correlated variables. Impurity-based importance is
biased toward continuous or high-cardinality features.

\section{Boosting as Functional Gradient Descent}

Boosting builds an additive model
\[
    F_M(x)=F_0(x)+\sum_{m=1}^M\nu\gamma_mh_m(x)
\]
by sequentially fitting weak learners. For differentiable empirical loss
$\sum_iL(y_i,F(x_i))$, iteration $m$ computes pseudo-residuals
\[
    r_{im}=-\left.
       \frac{\partial L(y_i,F)}{\partial F}
       \right|_{F=F_{m-1}(x_i)},
\]
fits $h_m$ to $(x_i,r_{im})$, chooses a step $\gamma_m$, and updates
$F_m=F_{m-1}+\nu\gamma_mh_m$. Thus boosting is gradient descent in function
space.

For squared loss, pseudo-residuals are ordinary residuals. For binary logistic
loss they are closely related to $y_i-\widehat p_i$. Tree depth controls
interaction order; learning rate $\nu$ and number of trees trade off. Smaller
$\nu$ usually requires more iterations and often generalizes better.

Boosting can overfit, especially with deep trees, noisy labels, or excessive
iterations. Tune on validation data and monitor training and validation loss.
Unlike bagging, trees are dependent and sequential: each corrects errors of
the current ensemble.

\begin{algorithm}[H]
\caption{Gradient boosting}
\begin{algorithmic}[1]
\State Initialize $F_0(x)=\argmin_c\sum_iL(y_i,c)$.
\For{$m=1,\ldots,M$}
  \State Compute $r_{im}=-\partial L(y_i,F)/\partial F$ at
  $F=F_{m-1}(x_i)$.
  \State Fit a weak learner $h_m$ to $\{(x_i,r_{im})\}$.
  \State Choose $\gamma_m$ by line search.
  \State Set $F_m=F_{m-1}+\nu\gamma_mh_m$.
\EndFor
\end{algorithmic}
\end{algorithm}
\section{Maximum Margin Classifiers}

For labels $y_i\in\{-1,1\}$, a linear classifier predicts
$\operatorname{sign}(\beta_0+x^\mathsf T\beta)$. For separable data, the
hard-margin SVM minimizes $\norm\beta^2/2$ subject to
$y_i(\beta_0+x_i^\mathsf T\beta)\geq1$. Its geometric margin is
$1/\norm\beta$. Feature scaling matters because distance and the penalty
depend on units.

\subsection{Soft-Margin SVM}

Nonseparable data use slack variables:
\begin{equation}
 \min_{\beta_0,\beta,\xi}\frac12\norm\beta^2+C\sum_i\xi_i
 \quad\text{subject to}\quad
 y_i(\beta_0+x_i^\mathsf T\beta)\geq1-\xi_i,\quad \xi_i\geq0.
 \label{eq:soft-margin-primal}
\end{equation}
Equivalently, minimize $\norm\beta^2/2+
C\sum_i[1-y_i(\beta_0+x_i^\mathsf T\beta)]_+$. Large $C$ gives less
regularization; small $C$ permits violations for a wider margin.

\subsection{SVM Duality}

Eliminating the primal variables gives
\begin{equation}
 \max_{\alpha}\left\{\sum_i\alpha_i-\frac12\sum_{i,j}
 \alpha_i\alpha_jy_iy_jx_i^\mathsf Tx_j\right\}
 \quad\text{s.t.}\quad
 0\leq\alpha_i\leq C,\quad \sum_i\alpha_iy_i=0.
 \label{eq:svm-dual}
\end{equation}
The primal coefficient is
$\widehat\beta=\sum_i\widehat\alpha_i y_ix_i$. Convexity and feasibility give
strong duality.

\subsection{Support Vectors and Kernels}

KKT complementary slackness shows that observations with $\alpha_i=0$ do not
directly affect the boundary. Those with $\alpha_i>0$ are support vectors:
$0<\alpha_i<C$ indicates a point on the margin, while $\alpha_i=C$ indicates
a point inside the margin or misclassified.

A positive-semidefinite kernel
$K(x,z)=\inner{\phi(x)}{\phi(z)}$ implicitly maps inputs to a feature space.
Replacing inner products in \eqref{eq:svm-dual} gives
\[
 f(x)=\widehat\beta_0+\sum_i\widehat\alpha_i y_iK(x_i,x).
\]
Polynomial and Gaussian radial kernels are common. Kernel width and $C$ must
be tuned jointly. The Gram matrix costs $O(n^2)$ memory. The raw margin is not
a probability; calibrate it on held-out data if probabilities are required.

\section{PCA from Constrained Optimization}

For centered $X\in\R^{n\times p}$, the first principal loading solves
\begin{equation}
 v_1=\argmax_{\norm v=1}\Var(Xv)
 =\argmax_{\norm v=1}v^\mathsf TX^\mathsf TXv.
 \label{eq:pca-rayleigh}
\end{equation}
The Lagrange equation makes $v_1$ the leading covariance eigenvector.
Subsequent loadings maximize variance subject to orthogonality. If
$X=UDV^\mathsf T$, scores are $UD$, loadings are $V$, and component $j$
explains variance proportional to $d_j^2$. Standardize columns when their
scales are not comparable. PCA is unsupervised: high variance need not predict
$Y$.

\subsection{PCA as LoRA}

Here LoRA means low-rank approximation. The Eckart--Young--Mirsky theorem
states
\begin{equation}
 X_k=U_kD_kV_k^\mathsf T
 =\argmin_{\rank(A)\leq k}\norm{X-A}_F^2,\qquad
 \norm{X-X_k}_F^2=\sum_{j>k}d_j^2.
 \label{eq:eckart-young}
\end{equation}
Modern low-rank adaptation uses an update $\Delta W=AB^\mathsf T$. It shares
the representation idea, but its factors optimize a supervised objective
rather than arise from a data SVD.

\subsection{PCR connection}

Principal components regression regresses $y$ on the first $k$ scores:
\[
 \widehat\beta_{\mathrm{PCR},k}=V_kD_k^{-1}U_k^\mathsf Ty.
\]
PCR discards low-variance directions; ridge shrinks them continuously. Because
PCR selects directions without $y$, a low-variance component can be highly
predictive. Select $k$ by CV rather than explained variance alone.

\section{K-means}

$K$-means minimizes
\begin{equation}
 \min_{C_1,\ldots,C_K,\,\mu_1,\ldots,\mu_K}
 \sum_{k=1}^K\sum_{i\in C_k}\norm{x_i-\mu_k}^2.
 \label{eq:kmeans-objective}
\end{equation}
Lloyd's algorithm alternates nearest-center assignment with updating each
center to its cluster mean. The objective decreases, but convergence is only
to a local optimum. Use $K$-means++ and multiple initializations, standardize
differently scaled features, and handle empty clusters explicitly.

$K$-means favors roughly spherical, similarly sized clusters and is sensitive
to outliers. Choose $K$ using stability, substantive usefulness, silhouette
summaries, or a model-based criterion; no single diagnostic defines the
scientifically correct number.

\section{Hierarchical Clustering}

Agglomerative clustering starts with singleton clusters and repeatedly merges
the closest pair. Common linkages are
\begin{align*}
 d_{\mathrm{single}}(A,B)&=\min_{i\in A,j\in B}d(i,j),\\
 d_{\mathrm{complete}}(A,B)&=\max_{i\in A,j\in B}d(i,j),\\
 d_{\mathrm{average}}(A,B)&=\frac1{\abs A\abs B}
     \sum_{i\in A,j\in B}d(i,j).
\end{align*}
Ward linkage chooses the smallest increase in within-cluster sum of squares.

The dendrogram records merge heights; branch order within a merge is not
meaningful. Single linkage can recover elongated clusters but chains through
nearby points. Complete linkage favors compact clusters. Distance, scaling,
and preprocessing define similarity. A dendrogram always displays structure,
even for homogeneous noise, so assess stability by perturbation or bootstrap
co-clustering and require substantive validation.




\end{document}
